% !TeX root = ../../Book4.tex
% 纲要标题：## 第19章　全导出函子、导出张量积与导出 Hom
% 纲要标签：b4:ch:18
% 纲要位置：计划纲要.md，第 3711 行。

\chapter{全导出函子、导出张量积与导出 Hom}
\label{b4:ch:18}
\BookChapterNumber{b4:ch:18}

\begin{chapterabstract}
本章以合适的复形替换构造全导出函子、导出张量积和导出 Hom，在明确有界条件与模的侧别后证明结合律、伴随和基变换，并讨论无界替换。
\end{chapterabstract}

\begin{learninggoals}
\begin{itemize}
\item 区分全导出函子、各次导出函子与超上同调，写出替换方向及泛性质中的比较态射。
\item 按右模与左模的侧别计算导出张量，并核对总次数及移位符号。
\item 用投射或内射替换计算导出 Hom，证明其与导出张量的伴随。
\item 分别检验基变换中的平坦性和有限生成性，识别普通基变换的失效。
\item 区分逐项投射、内射、平坦与相应的 K-性质。
\end{itemize}
\end{learninggoals}

若直接计算 \(\mathbb Z/2\mathbb Z\otimes_{\mathbb Z}\mathbb Z/2\mathbb Z\)，只能得到一个 \(\mathbb Z/2\mathbb Z\)。但将第一个模换成自由消解 \(0\rightarrow\mathbb Z\xrightarrow{2}\mathbb Z\rightarrow0\)，再与第二个模张量，乘 \(2\) 的微分变为零，所得复形在两个相邻次数各有一个 \(\mathbb Z/2\mathbb Z\) 的同调。直接张量丢掉的信息因此不是凭空添上的；它原本就在消解的另一度中。要让更换消解成为合法计算，还须说明何种复形可以替换、替换的选择是否改变结果，以及 Hom 的导出版本怎样与张量相配。

\ProseParagraphBreak
本章使用\chapref{b4:ch:17}中导出范畴的局部化和投射、内射模型，也使用\cref{b4:def:hom-complex}、\cref{b4:def:tensor-complex}的微分约定。与复合导出相关的谱序列已在\cref{b4:thm:grothendieck-spectral-sequence}建立。前三节的主要计算分别在有上界或有下界的范围进行；\subsecref{b4:subsec:2002:03}的无界余积与紧性讨论会调用\secref{b4:sec:1804}中的半自由替换。

\ProseParagraphBreak
对应阅读是 Weibel 的《An Introduction to Homological Algebra》\cite[Sections 10.5--10.8]{Weibel1994HomologicalAlgebra}：全导出构造见 Existence Theorem 10.5.6，张量与 Tor 见 Theorem 10.6.3 和 Corollary 10.6.4，导出 Hom 见 Theorem 10.7.4，复合条件见 Composition Theorem 10.8.2。本书在模范畴中展开基变换算例，并自行给出\secref{b4:sec:1804}所需的半自由替换构造。

% !TeX root = ../../Book4.tex
% 纲要标题：### 19.1 全导出函子
% 纲要标签：b4:sec:1801
% 纲要位置：计划纲要.md，第 3713 行。

\section{全导出函子}
\label{b4:sec:1801}

一个加性函子不一定保持拟同构，因而不能直接作用于导出态射。先选择能被它正确处理的复形模型，才能把导出群、自然变换与复合一起提升到导出范畴。

% 纲要内容（仅作写作计划，尚非教材正文）：
%
% * RF 和 LF 的构造及假设
% * 对象级导出函子与复形级导出的关系
% * 复合与自然变换
%

\subsection{\texorpdfstring{\(RF\) 与 \(LF\) 的构造条件}{RF 与 LF 的构造条件}}
\label{b4:subsec:1801:01}

对一个左正合函子逐次取导出，得到一列函子；但连接同态、自然变换及它们的复合仍需另外追踪。导出范畴允许保留产生这些群的整个复形。困难是原来的加性函子不一定保持拟同构，因而不能直接作用于分式表示。

\begin{definition}\label{b4:def:total-derived-functor}
设 \(\mathcal A,\mathcal B\) 是局部小交换范畴，\(F:\mathcal A\to\mathcal B\) 是加性函子，并固定导出范畴的集合大小约定。记 \(q_{\mathcal A}:K^+(\mathcal A)\to D^+(\mathcal A)\)、\(q_{\mathcal B}:K^+(\mathcal B)\to D^+(\mathcal B)\) 为局部化函子。三角函子
\(RF:D^+(\mathcal A)\to D^+(\mathcal B)\) 连同自然变换
\[
\eta:q_{\mathcal B}F\longrightarrow RFq_{\mathcal A}
\]
若满足下列泛性质，就称为 \(F\) 的\term[quan-you-daochu-hanzi]{全右导出函子}[total right derived functor]：对每个三角函子 \(T:D^+(\mathcal A)\to D^+(\mathcal B)\) 及自然变换 \(\xi:q_{\mathcal B}F\to Tq_{\mathcal A}\)，唯一存在自然变换 \(\theta:RF\to T\)，使 \(\xi=(\theta q_{\mathcal A})\eta\)。这里 \(F\) 在复形和同伦类上逐项作用。

对三角函子 \(LF:D^-(\mathcal A)\to D^-(\mathcal B)\)，将 \(q_{\mathcal A},q_{\mathcal B}\) 改为相应有上界范畴的局部化函子，并把自然变换方向改为
\(\epsilon:LFq_{\mathcal A}\to q_{\mathcal B}F\)。若对每个三角函子 \(T:D^-(\mathcal A)\to D^-(\mathcal B)\)，每个自然变换 \(Tq_{\mathcal A}\to q_{\mathcal B}F\) 都唯一经 \(T\to LF\) 分解，则称 \(LF\) 连同 \(\epsilon\) 为\term[quan-zuo-daochu-hanzi]{全左导出函子}[total left derived functor]。
\end{definition}
\symidx{\(RF,LF\)：全右导出函子与全左导出函子}

定义中的两个自然变换方向不同。内射替换给出 \(X\to I_X\)，因而右导出的比较来自 \(F(X)\to F(I_X)\)；投射替换给出 \(P_X\to X\)，因而左导出的比较来自 \(F(P_X)\to F(X)\)。加性函子未必保持这些拟同构，所以比较态射未必是同构。全导出函子修正的正是逐项 \(F\) 不能在导出范畴中直接使用的问题，不能要求修正前后的值总是相同。\cref{b4:prop:additive-right-derived-zero}将给出一个非零加性函子，其全右导出却为零。

\begin{theorem}\label{b4:thm:total-derived-functors}
设 \(\mathcal A,\mathcal B\) 是局部小交换范畴，\(F:\mathcal A\to\mathcal B\) 是加性函子，并在固定的集合大小约定下取导出范畴。
若 \(\mathcal A\) 有足够内射对象，则全右导出函子 \(RF:D^+(\mathcal A)\to D^+(\mathcal B)\) 存在，并对每个有下界上链复形 \(X\) 由
\[
RF(q_{\mathcal A}X)\cong q_{\mathcal B}F(I_X),\qquad X\xrightarrow{\simeq}I_X,
\]
计算，其中 \(I_X\) 是逐项内射的有下界复形。若 \(\mathcal A\) 有足够投射对象，则全左导出函子 \(LF:D^-(\mathcal A)\to D^-(\mathcal B)\) 存在，并对每个有上界上链复形 \(X\) 由
\[
LF(q_{\mathcal A}X)\cong q_{\mathcal B}F(P_X),\qquad P_X\xrightarrow{\simeq}X,
\]
计算，其中 \(P_X\) 是逐项投射的有上界复形。结果及其比较映射在导出范畴中不依赖替换的选择。
\end{theorem}
\begin{proof}
\cref{b4:thm:bounded-injective-model}给出的三角等价
\[
D^+(\mathcal A)\simeq K^+(\operatorname{Inj}\mathcal A)
\]
说明每个导出对象都有有下界的内射模型，而内射模型之间的导出态射可以按链同伦类计算。因此可先把导出对象换成内射模型，在模型上逐项作用 \(F\)，再进入 \(D^+(\mathcal B)\)。
记 \(\mathcal I=K^+(\operatorname{Inj}\mathcal A)\)，以
\(j:\mathcal I\to K^+(\mathcal A)\) 表示包含函子，并记此三角等价为
\(E=q_{\mathcal A}j:\mathcal I\to D^+(\mathcal A)\)。选择三角准逆
\(U:D^+(\mathcal A)\to\mathcal I\)，并将等价的自然同构取为伴随等价 \(U\dashv E\) 的单位与余单位：
\[
\rho:1_{D^+(\mathcal A)}\xrightarrow{\sim}EU,
\qquad
\kappa:UE\xrightarrow{\sim}1_{\mathcal I}.
\]
它们满足 \(E(\kappa_I)\rho_{E(I)}=1_{E(I)}\) 及
\(\kappa_{U(D)}U(\rho_D)=1_{U(D)}\)。定义
\(RF=q_{\mathcal B}FjU\)。加性函子保持有限双积、移位及锥的微分矩阵，所以逐项 \(F\) 是三角函子，所得复合也保持好三角。

对 \(X\in K^+(\mathcal A)\)，记 \(I_X=jU(q_{\mathcal A}X)\)。上一模型定理证明中的双射
\[
\operatorname{Hom}_{K^+(\mathcal A)}(X,I_X)
\xrightarrow{\sim}
\operatorname{Hom}_{D^+(\mathcal A)}(q_{\mathcal A}X,E U(q_{\mathcal A}X))
\]
将 \(\rho_{q_{\mathcal A}X}\) 唯一提升为同伦类 \(i_X:X\to I_X\)。因为其局部化是同构，\(i_X\) 的复形映射代表是拟同构。\(\rho\) 的自然性及上述双射的单射性，使 \(i\) 成为同伦范畴上的自然变换。这并不要求预先自然地选定严格链映射代表：更换代表只会改变一个链同伦，而加性 \(F\) 保持链同伦。因此
\(\eta_X=q_{\mathcal B}F(i_X)\) 不依赖代表，给出所需自然比较。

给定 \(T,\xi\)，先看泛分解等式怎样决定所需变换。第一条三角恒等式及提升唯一性给出 \(i_{jI}=j(\kappa_I^{-1})\)，第二条三角恒等式因而说明
\[
RF(\rho_D)=q_{\mathcal B}Fj(U(\rho_D))
=q_{\mathcal B}Fj(\kappa_{U(D)}^{-1})
=\eta_{jU(D)}.
\]
若 \(\theta:RF\to T\) 满足泛分解等式，沿 \(\rho_D:D\xrightarrow{\sim}EU(D)\) 的自然性便要求
\[
T(\rho_D)\theta_D
=\theta_{EU(D)}RF(\rho_D)
=\theta_{EU(D)}\eta_{jU(D)}
=\xi_{jU(D)}.
\]
由于 \(T(\rho_D)\) 可逆，\(\theta_D\) 只能取为
\[
\theta_D=(T(\rho_D))^{-1}\xi_{jU(D)}:
RF(D)\longrightarrow T(D).
\]
这里 \(\xi_{jU(D)}\) 的目标是 \(T(EU(D))\)；它与 \(T(\rho_D)\) 的逆可以复合。\(\xi\) 与 \(\rho\) 的自然性说明 \(\theta\) 自然，而
\[
\theta_{q_{\mathcal A}X}\eta_X
=(T(\rho_{q_{\mathcal A}X}))^{-1}
T(q_{\mathcal A}i_X)\xi_X=\xi_X.
\]
这证明存在性；前面的强制公式同时证明唯一性。

左导出采用 \(P_X\to X\) 的投射替换，所以比较方向是 \(LF\to F\)。给定的变换是 \(\xi:Tq_{\mathcal A}\to q_{\mathcal B}F\)，要找的是 \(\theta:T\to LF\)，使 \(\epsilon_X\theta_{q_{\mathcal A}X}=\xi_X\)；它的泛分解方向与右导出相反。记
\(\mathcal P=K^-(\operatorname{Proj}\mathcal A)\)，令
\(j^-:\mathcal P\to K^-(\mathcal A)\) 为包含，
\(E^-=q_{\mathcal A}j^-\)。由\cref{b4:thm:bounded-projective-model}选择三角准逆
\(V:D^-(\mathcal A)\to\mathcal P\)，以及伴随等价 \(E^-\dashv V\) 的单位、余单位
\[
\nu:1_{\mathcal P}\xrightarrow{\sim}VE^-,
\qquad \lambda:E^-V\xrightarrow{\sim}1_{D^-(\mathcal A)}.
\]
投射模型的 Hom 双射将 \(\lambda_{q_{\mathcal A}X}\) 唯一提升为自然的同伦类
\(p_X:j^-V(q_{\mathcal A}X)\to X\)。定义
\(LF=q_{\mathcal B}Fj^-V\)、\(\epsilon_X=q_{\mathcal B}F(p_X)\)。三角恒等式给出
\(LF(\lambda_D)=\epsilon_{j^-V(D)}\)。沿 \(\lambda_D:E^-V(D)\xrightarrow{\sim}D\) 的自然性和泛分解等式因而强制
\[
\theta_D T(\lambda_D)
=LF(\lambda_D)\theta_{E^-V(D)}
=\epsilon_{j^-V(D)}\theta_{E^-V(D)}
=\xi_{j^-V(D)}.
\]
所以所需变换在导出对象 \(D\) 上为
\[
\theta_D=\xi_{j^-V(D)}(T(\lambda_D))^{-1}:
T(D)\longrightarrow LF(D).
\]
自然性及 \(\epsilon_X\theta_{q_{\mathcal A}X}=\xi_X\) 由 \(\xi\) 与 \(\lambda\) 的自然性方块得到；前面的强制公式证明唯一性。

任意另一内射或投射替换都经模型的 Hom 双射得到与给定替换相容的唯一同伦比较；它在局部化后为同构，因模型等价全忠实而是同伦等价。逐项 \(F\) 保持同伦等价，故这些替换计算的对象典范同构于上述函子的值。泛性质也使不同模型选择所得函子连同比较变换唯一自然同构。
\end{proof}

这里构造存在只需加性，不要求左正合或右正合。正合性将用于把零次同调重新识别为 \(F\) 本身。以上泛性质及构造可对照 Weibel\cite[Definition 10.5.1, Existence Theorem 10.5.6]{Weibel1994HomologicalAlgebra}。

\subsection{对象级导出与复形级导出}
\label{b4:subsec:1801:02}

把对象 \(A\in\mathcal A\) 看作零次复形，替换 \(A\to I^\bullet\) 就是通常的内射消解。因此，当 \(F\) 左正合且 \(\mathcal A\) 有足够内射对象时，
\[
H^n(RF(A))\simeq R^nF(A)\quad(n\ge0),\qquad
H^n(RF(A))=0\quad(n<0).
\]
尤其 \(H^0(RF(A))=F(A)\)。右正合函子在有足够投射对象时给出
\(H^{-n}(LF(A))=L_nF(A)\)。负号来自本书始终使用升次微分，而通常的投射消解用非负链次数编号。

\ProseParagraphBreak
若 \(F\) 左正合、\(\mathcal A\) 有足够内射对象，而 \(X\) 为有下界复形，则 \(H^n(RF(X))\) 就是\cref{b4:def:hypercohomology}中的 \(\mathbb H^n(F,X)\)：两种计算选择同一内射复形后相同。\cref{b4:thm:hypercohomology-spectral-sequences}的第二套谱序列于是写成
\[
E_2^{p,q}=R^pF(H^q(X))
\Longrightarrow H^{p+q}(RF(X)).
\]
这里 \(p\ge0\)，\(q\) 有下界，每个总次数的滤过有限。即使知道所有第二页数据，仍须经过后续微分，并由终页数据得到目标的逐层商，不能一般地把它们直接相加。\(R^nF(A)\) 是单个对象的第 \(n\) 次导出，\(RF(X)\) 是整个复形的导出对象，\(H^n(RF(X))\) 才是该对象的第 \(n\) 次上同调。特别地，\(RF(X)\) 不能只由 \(R^nF(H^0X)\) 代替；即使 \(F\) 是恒等函子，\(RF(X)=X\) 仍保留每个 \(H^nX\)。

下文沿用\cref{b4:def:functor-acyclic}的“\(F\)-零调对象”：当 \(F\) 左正合时，它只要求 \(R^jF(Q)=0\) 对 \(j>0\) 成立；“无上同调对象”是同一性质的别称。它与整个复形的零调性是两个不同条件。

\begin{proposition}\label{b4:prop:total-derived-acyclic-complex}
设 \(F:\mathcal A\to\mathcal B\) 是局部小交换范畴间的加性左正合函子，\(\mathcal A\) 有足够内射对象。若 \(X\) 是 \(\mathcal A\) 中的有下界上链复形，且每项均满足
\(R^jF(X^n)=0\) 对所有 \(j>0\) 成立，则自然映射 \(F(X)\to RF(X)\) 是拟同构。若 \(F\) 为加性右正合函子，\(\mathcal A\) 有足够投射对象，\(X\) 为 \(\mathcal A\) 中有上界的上链复形，并且 \(L_jF(X^n)=0\) 对所有 \(n\in\mathbb Z\)、\(j>0\) 成立，则对偶的自然映射 \(LF(X)\to F(X)\) 是拟同构。
\end{proposition}
\begin{proof}
先设 \(E\) 是有下界的正合复形，每项为 \(F\)-零调对象。记 \(Z^n=\ker \mathrm{d}_E^n\)。从低于所有非零项的次数开始，\(Z^n=0\)。短正合列
\[
0\longrightarrow Z^n\longrightarrow E^n\longrightarrow Z^{n+1}\longrightarrow0
\]
给出导出长正合列。若 \(Z^n\) 已 \(F\)-零调，其中的片段
\[
F(E^n)\longrightarrow F(Z^{n+1})\longrightarrow R^1F(Z^n)=0
\]
使第一个映射满；而 \(E^n\) 的零调性给出
\[
R^iF(Z^{n+1})\cong R^{i+1}F(Z^n)=0\qquad(i\ge1),
\]
使下一层 \(Z^{n+1}\) 也 \(F\)-零调。左正合性再给出
\(0\to FZ^n\to FE^n\to FZ^{n+1}\to0\) 正合。由起始的 \(Z^n=0\) 逐次向上归纳，\(F(E)\) 便正合。

取 \(X\to I\) 的有下界内射替换，其锥正合、有下界，且每项是内射对象与某个 \(X^n\) 的有限直和，故逐项 \(F\)-零调。上段说明 \(F(\operatorname{Cone}(X\to I))\) 正合。加性保证此复形是 \(\operatorname{Cone}(FX\to FI)\)，锥判据得到拟同构。对偶情形从最高次数向下归纳。
\end{proof}

上一个证明说明了有下界假设的实际用途：它给余圈对象的归纳提供起点。逐项对 \(F\) 零调是对象性质，整个复形经 \(F\) 后仍正合是另一个性质；无界时不能仅凭前者推出后者。

\subsection{复合与自然变换}
\label{b4:subsec:1801:03}

若 \(\alpha:F\to G\) 是两个加性函子间的自然变换，在同一个内射替换上逐项取 \(\alpha_{I_X}\)，便得到
\(R\alpha:RF\to RG\)。比较态射使它自然，且
\(R(\beta\alpha)=(R\beta)(R\alpha)\)、\(R(1_F)=1_{RF}\)。左导出使用同一个投射替换，也有相同结论。因而自然变换不是在每个导出群上重新猜测，而是先在复形层面构造。

\begin{theorem}\label{b4:thm:total-derived-composition}
设 \(\mathcal A,\mathcal B,\mathcal C\) 是局部小交换范畴，\(\mathcal A,\mathcal B\) 有足够内射对象，\(F:\mathcal A\to\mathcal B\)、\(G:\mathcal B\to\mathcal C\) 为加性左正合函子。总有自然比较
\[
R(GF)\longrightarrow RG\circ RF
\quad\text{在}\;D^+(\mathcal A)\;\text{上}.
\]
若 \(F\) 把内射对象送到 \(G\)-零调对象，则这个比较是自然同构。
\end{theorem}
\begin{proof}
取 \(X\to I\)，再取 \(F(I)\to J\) 的有下界内射替换。\(R(GF)(X)\) 由 \(GF(I)\) 表示，\(RG(RF(X))\) 由 \(G(J)\) 表示，比较就是
\(G(F(I))\to G(J)\)：第一次导出使用 \(I\)，第二次导出还要把 \(F(I)\) 换成内射模型 \(J\)。同伦比较保证选择独立与自然性；等价地，它由全右导出的泛性质作用于两次比较变换的复合产生。虽然 \(F(I)\to J\) 为拟同构，加性左正合的 \(G\) 未必保持它，所以这个比较并不自动可逆。

若每个 \(F(I^n)\) 都 \(G\)-零调，则\cref{b4:prop:total-derived-acyclic-complex}应用于有下界复形 \(F(I)\)，说明上述比较是拟同构。这也给出两种复合方式之间的自然同构。
\end{proof}

若 \(F\) 把内射对象送到 \(G\)-零调对象，则对一般
\(X\in D^+(\mathcal A)\)，将\cref{b4:thm:hypercohomology-spectral-sequences}的第二套谱序列应用于 \(RF(X)\) 的有下界复形模型，得到自然的强收敛谱序列
\[
E_2^{p,q}=R^pG\bigl(H^q(RF(X))\bigr)
\Longrightarrow H^{p+q}\bigl(R(GF)(X)\bigr).
\]
这里 \(p\ge0\)，\(q\) 有下界，所以每个总次数的滤过仍有限。若特别取
\(X=A[0]\)，其中 \(A\in\mathcal A\)，则
\(H^q(RF(A[0]))=R^qF(A)\) 对 \(q\ge0\) 成立，负次数为零。于是恢复\cref{b4:thm:grothendieck-spectral-sequence}的第一象限谱序列
\[
E_2^{p,q}=R^pG(R^qF(A))
\Longrightarrow R^{p+q}(GF)(A).
\]
它提供目标的滤过及其逐层商，而不自动给出直和分解
\[
R^n(GF)(A)\cong\bigoplus_{p+q=n}R^pG(R^qF(A)).
\]
全导出复合的比较与零调性条件见 Weibel\cite[Composition Theorem 10.8.2, Corollary 10.8.3]{Weibel1994HomologicalAlgebra}。

\ProseParagraphBreak
对右正合函子及投射替换有对偶版本：若 \(F\) 把投射对象送到 \(G\)-零调对象，则 \(L(GF)\simeq LG\circ LF\)；一般比较的方向为 \(LG\circ LF\to L(GF)\)。如果不检查零调性，就没有理由让比较成为同构。

\begin{proposition}\label{b4:prop:additive-right-derived-zero}
在交换群范畴 \(\mathbf{Ab}\) 上取加性函子
\[
F(M)=M\otimes_{\mathbb Z}\mathbb Z/2\mathbb Z.
\]
其全右导出函子 \(RF:D^+(\mathbf{Ab})\to D^+(\mathbf{Ab})\) 为零函子，但 \(F(\mathbb Z)=\mathbb Z/2\mathbb Z\ne0\)。因此在零次复形 \(\mathbb Z[0]\) 上，比较态射 \(F(\mathbb Z)[0]\to RF(\mathbb Z[0])\) 不是同构。仅有加性不能保证对每个交换群 \(M\) 都有 \(H^0(RF(M[0]))\cong F(M)\)。
\end{proposition}
\begin{proof}
设 \(I\) 是内射交换群。对任意 \(x\in I\)，将同态 \(2\mathbb Z\to I\)、\(2\mapsto x\) 沿包含 \(2\mathbb Z\hookrightarrow\mathbb Z\) 延拓为 \(f:\mathbb Z\to I\)。令 \(y=f(1)\)，便有 \(2y=f(2)=x\)。故乘 \(2\) 在 \(I\) 上满射，\(I/2I=0\)，从而
\[
F(I)\cong I/2I=0.
\]
对任意有下界复形 \(X\)，取有下界的逐项内射替换 \(X\to I_X\)。\(F(I_X)\) 逐项为零，所以\cref{b4:thm:total-derived-functors}给出的 \(RF\) 是零函子。另一方面，\(F(\mathbb Z)=\mathbb Z/2\mathbb Z\) 非零，其零次复形的 \(H^0\) 非零，不能同构于零导出对象 \(RF(\mathbb Z[0])\)。该函子为右正合函子，但它把单射 \(\mathbb Z\xrightarrow{2}\mathbb Z\) 送为 \(\mathbb Z/2\mathbb Z\) 上的零映射，故不是左正合函子。因此它的全右导出不受前面左正合情形的零次识别约束。
\end{proof}

% !TeX root = ../../Book4.tex
% 纲要标题：### 19.2 导出张量积
% 纲要标签：b4:sec:1802
% 纲要位置：计划纲要.md，第 3719 行。

\section{导出张量积}
\label{b4:sec:1802}

普通张量积可能在更换拟同构模型以后改变同调。平坦或投射替换修正这一问题，并使张量积的结合律、单位同构及 Tor 计算相互一致。有上界平坦复形的论证使用\cref{b4:prop:filtered-colimits-cohomology}中滤过余极限与同调交换的性质。

% 纲要内容（仅作写作计划，尚非教材正文）：
%
% * 平坦或投射替换
% * 结合律与适当的单位同构
% * Tor 作为导出张量的同调
%

\subsection{平坦替换与投射替换}
\label{b4:subsec:1802:01}

普通张量积不能保持所有拟同构。例如整数模复形 \(\mathbb Z\xrightarrow{2}\mathbb Z\)（次数 \(-1,0\)）拟同构于 \(\mathbb Z/2\mathbb Z\)，但与 \(\mathbb Z/2\mathbb Z\) 张量后前者多出一次负次数同调。为了让张量积在导出范畴中有意义，需要能保存拟同构的替换。

\begin{lemma}\label{b4:lem:bounded-flat-tensor-preserves}
设 \(R\) 为环，\(P\) 为逐项平坦且有上界的右 \(R\)-模复形，\(A\) 为任意正合的左 \(R\)-模复形。采用直和总化时，\(P\otimes_RA\) 正合。因此 \(P\otimes_R-\) 保持任意复形之间的拟同构。交换左右侧别后也成立。
\end{lemma}
\begin{proof}
若 \(P\) 只有一项，平坦性说明张量保留 \(A\) 的各个核像短列，所以结果正合。若 \(P\) 有有限个非零项，按最高非零项逐次剥去，所得复形短列在分次模意义下分裂；张量后仍是短正合列，用同调长正合列归纳，得到正合性。

一般情形取粗截断 \(P_{\ge m}\)，即低于 \(m\) 的项换成零。它是 \(P\) 的子复形，且因有上界而只有有限个非零项。截口可能改变同调，故这些粗截断一般不拟同构于 \(P\)；这里使用的是随着 \(m\) 向负方向移动，\(P\) 为它们的逐项滤过余极限。张量与直和总化都与此余极限交换，而模范畴的滤过余极限正合并与同调交换，见\cref{b4:prop:filtered-colimits-cohomology}。因此
\[
 H^n(P\otimes_RA)
 \cong\varinjlim_{m\to-\infty}H^n(P_{\ge m}\otimes_RA)=0.
\]

若 \(u:A\to B\) 是拟同构，其锥正合。张量锥与张量映射的锥由下面的分量识别自然同构：
\[
 \begin{gathered}
 \theta:P\otimes_R\operatorname{Cone}(u)
 \longrightarrow\operatorname{Cone}(1_P\otimes u),\\
 x\otimes(b,a)\longmapsto
 \bigl(x\otimes b,(-1)^p x\otimes a\bigr),
 \end{gathered}
\]
其中 \(x\in P^p\)、\(b\in B^q\)、\(a\in A^{q+1}=A[1]^q\)。第二分量的符号区别了先移位右因子与先张量再移位。具体地，在 \(x\otimes(0,a)\) 上，两边的微分计算均给出
\[
 \begin{aligned}
 \theta\mathrm{d}\bigl(x\otimes(0,a)\bigr)
 &=\mathrm{d}\theta\bigl(x\otimes(0,a)\bigr)\\
 &=\Bigl((-1)^p x\otimes u(a),\;
 (-1)^{p+1}\cdot\mathrm{d}_Px\otimes a-x\otimes\mathrm{d}_Aa\Bigr).
 \end{aligned}
\]
在 \(x\otimes(b,0)\) 上，\(\theta\) 不改变分量，微分相容性直接由张量微分得到。\(\theta\) 的逆使用同一符号。于是张量后的锥正合，得到后一结论。
\end{proof}

\begin{theorem}\label{b4:thm:derived-tensor}
设 \(R\) 是环。对有上界的右 \(R\)-模复形 \(X\) 和有上界的左 \(R\)-模复形 \(Y\)，选择逐项投射的有上界替换 \(p:P\to X\)、\(q:Q\to Y\)。公式
\[
X\otimes_R^{\mathbf L}Y\coloneqq P\otimes_RQ
\]
定义双函子
\[
-\otimes_R^{\mathbf L}-:
D^-(\operatorname{Mod}\text{-}R)\times D^-(R\text{-}\operatorname{Mod})
\longrightarrow D^-(\mathbf{Ab}).
\]
它称为\term[daochu-zhangliangji]{导出张量积}[derived tensor product]。也可以只替换一个变量，使用 \(P\otimes_RY\) 或 \(X\otimes_RQ\)。若 \(P'\to X\) 或 \(Q'\to Y\) 是有上界的逐项平坦替换，则 \(P'\otimes_RY\) 或 \(X\otimes_RQ'\) 分别典范同构于这个双函子的值；同时使用两侧的平坦替换也得到同一导出对象。
\end{theorem}
\symidx{\(\operatorname{Mod}\text{-}R\)：幺性右模范畴}
\begin{proof}
上一引理分别应用于 \(P\) 和 \(Q\)，说明自然映射
\[
P\otimes_RY\xleftarrow{\,1_P\otimes q\,}P\otimes_RQ
\xrightarrow{\,p\otimes1_Q\,}X\otimes_RQ
\]
都是拟同构。左箭头使用 \(P\) 有上界且逐项投射，从而与 \(P\) 张量保持 \(q\) 的拟同构；右箭头使用 \(Q\) 的对应性质，从而与 \(Q\) 张量保持 \(p\) 的拟同构。因此固定一侧的适当模型，另一侧就可以保留原复形。投射模型中的态射由同伦类表示，由\cref{b4:prop:tensor-hom-preserve-homotopy}，张量积在两个变量都保持同伦，所以先在两个投射同伦范畴上得到双函子，再由两侧的投射模型等价下降到所写导出范畴。

若 \(p':P'\to X\) 是有上界逐项平坦替换，由\cref{b4:thm:bounded-projective-model}证明中的比较结论，存在唯一的同伦类 \(t:P\to P'\)，使 \(p't=p\) 于同伦范畴中。由于 \(p,p'\) 都是拟同构，\(t\) 也是拟同构。于是
\[
P\otimes_RQ
\xrightarrow{\,t\otimes1_Q\,}P'\otimes_RQ
\xrightarrow{\,1_{P'}\otimes q\,}P'\otimes_RY
\]
的两条箭头都是拟同构：第一条使用 \(Q\) 的平坦性及引理的左右对偶，第二条使用 \(P'\) 的平坦性。张量保持同伦，所以这个合成不依赖 \(t\) 的代表元。若更换 \(P\) 或 \(Q\)，投射模型的比较同伦类分别满足与 \(p\) 或 \(q\) 的相容等式；\(t\) 的唯一性使上述合成与模型间的比较交换。因此它给出 \(P'\otimes_RY\) 与已定义双函子值的典范同构。交换左右侧别得到第二变量的结论，再依次替换两侧即可同时使用 \(P',Q'\)。双函子在态射上的作用始终由投射模型定义。两变量有上界使每个总次数只涉及有限条对角项，且输出仍有上界。
\end{proof}
\symidx{\(X\otimes_R^{\mathbf L}Y\)：导出张量积}

对于非交换环，上式一般只给交换群复形；不能在结果上无说明添一个左 \(R\)-模结构。若 \(X\) 还带与右 \(R\)-作用交换的左 \(S\)-作用，则张量复形上的作用为 \(s\cdot(x\otimes y)=(sx)\otimes y\)。可以固定 \(X\)，仅用左 \(R\)-模投射复形 \(Q\) 替换 \(Y\)，使 \(X\otimes_RQ\) 按这个公式保留左 \(S\)-作用，且微分仍为 \(S\)-线性的。若仅将 \(X\) 看作右 \(R\)-模而任取替换 \(P\to X\)，\(P\) 未必带有与微分相容的左 \(S\)-作用，此时 \(P\otimes_RY\) 也不能按上述公式成为左 \(S\)-模复形。若要替换 \(X\)，就应说明替换如何保留所需双模结构。模的有上界导出张量积可对照 Weibel\cite[Definition 10.6.1, Lemma 10.6.2, Theorem 10.6.3]{Weibel1994HomologicalAlgebra}。

\subsection{结合律与单位同构}
\label{b4:subsec:1802:02}

\begin{proposition}\label{b4:prop:derived-tensor-associativity}
设 \(R\) 为交换环，\(X,Y,Z\in D^-(R\text{-}\operatorname{Mod})\)。导出张量积有自然同构
\[
(X\otimes_R^{\mathbf L}Y)\otimes_R^{\mathbf L}Z
\simeq X\otimes_R^{\mathbf L}(Y\otimes_R^{\mathbf L}Z),
\qquad
R\otimes_R^{\mathbf L}X\simeq X\simeq X\otimes_R^{\mathbf L}R.
\]
还有对称同构 \(X\otimes_R^{\mathbf L}Y\simeq Y\otimes_R^{\mathbf L}X\)。若 \(P,Q\) 分别是 \(X,Y\) 的有上界逐项平坦替换，它在齐次元上由 \(x\otimes y\mapsto(-1)^{pq}y\otimes x\) 表示，其中 \(x\in P^p\)、\(y\in Q^q\)。
\end{proposition}
\begin{proof}
取有上界逐项投射替换 \(P,Q,U\)。交换环上两个投射模的张量积仍投射：把二者写成自由模的直和项，其张量就是自由模张量的直和项。每条总次数对角线有限，故 \(P\otimes_RQ\) 及 \(Q\otimes_RU\) 均逐项投射、有上界，可以直接用来计算下一次导出张量。

两边都由三重张量总复形表示，齐次纯张量的微分为
\[
\mathrm{d}(p\otimes q\otimes u)=\mathrm{d}p\otimes q\otimes u
+(-1)^{|p|}p\otimes \mathrm{d}q\otimes u
+(-1)^{|p|+|q|}p\otimes q\otimes \mathrm{d}u.
\]
普通张量的结合映射保持此公式，给出复形同构；对四个因子，所有结合映射都只改变括号，所以它们的不同复合相同。\(R[0]\) 本身投射，通常的单位映射逐项给出单位同构。交换两个因子的公式中，微分越过次数 \(p\) 或 \(q\) 时产生的符号正好补偿 \((-1)^{pq}\) 的改变，直接代入便得到复形映射；两次交换的符号是 \((-1)^{2pq}=1\)，故为同构。这些复形映射与投射模型中的比较态射相容，故在导出范畴中仍自然。
\end{proof}

单位结论不需要交换性：对任意环 \(R\)，右模复形 \(X\) 满足 \(X\otimes_R^{\mathbf L}R\simeq X\)，左模复形 \(Y\) 满足 \(R\otimes_R^{\mathbf L}Y\simeq Y\)，此处 \(R\) 带通常的 \((R,R)\)-双模结构。非交换环上的多重张量也有结合映射，但每两个相邻因子必须具有匹配的左右作用，并须核对保留这些作用的替换；上面的交换环命题没有免除这些要求。

\subsection{Tor 与导出张量的同调}
\label{b4:subsec:1802:03}

\begin{corollary}\label{b4:cor:tor-derived-tensor}
设 \(R\) 为环，\(M\) 为右 \(R\)-模，\(N\) 为左 \(R\)-模，二者视作零次复形。则对每个 \(n\ge0\)，有自然同构
\[
H^{-n}(M\otimes_R^{\mathbf L}N)\simeq\operatorname{Tor}_n^R(M,N),
\]
而正次数同调为零。特别地，零次同调为 \(M\otimes_RN\)。
\end{corollary}
\begin{proof}
把 \(M\) 的投射链消解 \(P_\bullet\to M\) 重编号为 \(P^{-n}=P_n\)，并只替换第一变量。张量复形在 \(-n\) 次的同调正是\cref{b4:thm:tor-balanced}所计算的 \(\operatorname{Tor}_n\)。其复形没有正次数项，零次同调由张量的右正合性识别。
\end{proof}

例如，设 \(R\) 为交换环，\(a\in R\) 不是零因子。\(R/(a)\) 由 \(R\xrightarrow{a}R\) 的两项自由复形替换，因此对任意 \(R\)-模 \(N\)，
\[
R/(a)\otimes_R^{\mathbf L}N\simeq
\bigl(N\xrightarrow{a}N\bigr),\qquad
H^{-1}=\{\,n\in N\mid an=0\,\},\quad H^0=N/aN.
\]
若 \(a\) 同时零化 \(N\)，微分为零，故在导出范畴中这是 \(N\oplus N[1]\)。若 \(a\) 在 \(N\) 上仍为非零因子，则负次同调消失，此时才可用普通张量积代替。\(a\) 在 \(R\) 上不是零因子的条件保证所用两项自由复形确为消解。

\ProseParagraphBreak
把这个公式用于两条直线，便能看出负次同调记录了什么。设 \(k\) 为域，\(R=k[x,y]\)，并取 \(M=R/(x)\)。若 \(N=R/(y)=k[x]\)，则乘 \(x\) 单射，因而
\[
 H^{-1}(M\otimes_R^{\mathbf L}N)=0,\qquad
 H^0(M\otimes_R^{\mathbf L}N)=k.
\]
两条坐标轴的方程彼此独立，普通张量积已经给出它们的交点。若改取 \(N=M=k[y]\)，则乘 \(x\) 为零，故
\[
 M\otimes_R^{\mathbf L}M\simeq M\oplus M[1],\qquad
 H^{-1}=M,
\]
而普通张量积只有 \(H^0=M\)。这一次两个模使用同一条方程，新增的负次同调保留了方程重复所产生的关系。导出张量积既记录共同的商，也记录形成这个商时关系之间的依赖。

% !TeX root = ../../Book4.tex
% 纲要标题：### 19.3 导出 Hom
% 纲要标签：b4:sec:1803
% 纲要位置：计划纲要.md，第 2964 行。

\section{导出 Hom}
\label{b4:sec:1803}

导出 Hom 保留全部 Ext 次数，导出 Hom--张量伴随则联系两个不同的计算模型。基变换需要额外假设，双数代数的周期消解给出非平坦情形的明确反例。

% 纲要内容（仅作写作计划，尚非教材正文）：
%
% * RHom 与 Ext
% * 导出 Hom–张量伴随
% * 基变换的充分条件及失败例子
% * 以 \(R=k[\varepsilon]/(\varepsilon^2)\) 和 \(k=R/(\varepsilon)\) 对照普通张量与导出张量：周期自由消解给出非零的高阶 Tor，说明非平坦基变换丢失的信息
%

\subsection{RHom 与 Ext}
\label{b4:subsec:1803:01}

\cref{b4:def:hom-complex}已经定义
\[
\operatorname{Hom}_R^n(X,Y)=
\prod_p\operatorname{Hom}_R(X^p,Y^{p+n}),\qquad
\delta f=\mathrm{d}_Yf-(-1)^nf\mathrm{d}_X.
\]
零次余圈是复形映射，零次余边是零同伦映射。更一般地，次数 \(n\) 的余圈满足 \(\mathrm{d}_Yf=(-1)^nf\mathrm{d}_X\)，正好给出 \(X\to Y[n]\) 的复形映射，因此
\(H^n\operatorname{Hom}_R^\bullet(X,Y)=\operatorname{Hom}_{K(R)}(X,Y[n])\)。把同伦范畴换成导出范畴，也需要修正 Hom 复形。

\begin{theorem}\label{b4:thm:derived-hom}
设 \(R\) 为环，记 \(D(R)=D(R\text{-}\mathbf{Mod})\)。设 \(X\) 为有上界的左 \(R\)-模上链复形，\(Y\) 为有下界的左 \(R\)-模上链复形。取有上界投射替换 \(P\to X\) 与有下界内射替换 \(Y\to I\)。三个交换群复形
\[
\operatorname{Hom}_R^\bullet(P,Y),\qquad
\operatorname{Hom}_R^\bullet(P,I),\qquad
\operatorname{Hom}_R^\bullet(X,I)
\]
自然拟同构，定义
\[
\mathbf R\operatorname{Hom}_R:
D^-(R\text{-}\operatorname{Mod})^{\mathrm{op}}
\times D^+(R\text{-}\operatorname{Mod})
\longrightarrow D^+(\mathbf{Ab}).
\]
这称为\term[daochu-Hom]{导出 Hom}[derived Hom]。其同调满足
\[
H^n\mathbf R\operatorname{Hom}_R(X,Y)
\simeq\operatorname{Hom}_{D(R)}(X,Y[n])
\quad(n\in\mathbb Z).
\]
若 \(R\) 交换，所构造的 Hom 复形带自然的 \(R\)-模结构，输出可取在 \(D^+(R\text{-}\operatorname{Mod})\) 中。
\end{theorem}
\begin{proof}
\cref{b4:thm:bounded-projective-model}的比较结论说明，对于任意正合复形 \(A\)，所有 \(P\to A[n]\) 都零同伦，所以 \(\operatorname{Hom}_R^\bullet(P,A)\) 正合。\cref{b4:thm:bounded-injective-model}对偶地说明 \(\operatorname{Hom}_R^\bullet(A,I)\) 正合。因此对 \(Y\to I\) 和 \(P\to X\) 的正合锥施加 Hom，便得到所写的两条拟同构。Hom 与锥相容时使用其定义中的符号，或在每次同调上直接用上述同伦类解释，都会给出同一结论。

投射及内射比较态射在同伦意义下唯一，因此此构造在两个变量上自然，并把拟同构送到同构，从而定义所写双函子。若 \(P^p=0\) 对 \(p>a\)，\(I^j=0\) 对 \(j<b\)，则 \(\operatorname{Hom}^n(P,I)=0\) 对 \(n<b-a\)，所以输出有下界。最后
\[
H^n\operatorname{Hom}_R^\bullet(P,I)
=\operatorname{Hom}_{K(R)}(P,I[n])
\simeq\operatorname{Hom}_{D(R)}(X,Y[n]),
\]
最后一步仍是投射或内射比较结论。交换环时，标量与所有映射及比较相容，故得到 \(R\)-模版本。
\end{proof}
\symidx{\(\mathbf R\operatorname{Hom}_R(X,Y)\)：导出 Hom}

对模 \(M,N\) 置于零次，以上公式化为
\[
H^n\mathbf R\operatorname{Hom}_R(M,N)=\operatorname{Ext}_R^n(M,N)
\quad(n\ge0),
\]
负次为零，与\cref{b4:thm:derived-ext-identification}一致。把所有 Ext 群单独列出，通常不能恢复导出 Hom 复形；后者还保留它在三角和复合中的相容关系。只替换第二变量时可以允许第一变量无界，只替换第一变量时可以允许第二变量无界；本节的共同有界范围使两种计算可以直接比较。参阅 Weibel\cite[Definition 10.7.2, Theorem 10.7.4, Corollary 10.7.5]{Weibel1994HomologicalAlgebra}。

\subsection{导出 Hom–张量伴随}
\label{b4:subsec:1803:02}

普通的 Hom--张量伴随把双线性映射写成一族线性映射。复形版本仍用同一个代入公式，不过必须核对各个次数和微分。

\begin{theorem}\label{b4:thm:derived-hom-tensor-adjunction}
设 \(R\) 为交换环，记 \(D(R)=D(R\text{-}\mathbf{Mod})\)。设 \(X,Y\in D^-(R\text{-}\operatorname{Mod})\)，\(Z\in D^+(R\text{-}\operatorname{Mod})\)。存在自然同构
\[
\mathbf R\operatorname{Hom}_R(X\otimes_R^{\mathbf L}Y,Z)
\simeq
\mathbf R\operatorname{Hom}_R
\bigl(X,\mathbf R\operatorname{Hom}_R(Y,Z)\bigr).
\]
特别地，在零次同调上得到
\[
\operatorname{Hom}_{D(R)}(X\otimes_R^{\mathbf L}Y,Z)
\simeq
\operatorname{Hom}_{D(R)}
\bigl(X,\mathbf R\operatorname{Hom}_R(Y,Z)\bigr).
\]
\end{theorem}
\begin{proof}
取有上界投射替换 \(P\to X\)、\(Q\to Y\)，以及有下界内射替换 \(Z\to I\)。\(P\otimes_RQ\) 是有上界逐项投射复形：交换环上两个投射模的张量是投射模，因为它是两个自由模张量的直和项。故左侧由 \(\operatorname{Hom}^\bullet_R(P\otimes_RQ,I)\) 表示。

内层 \(\mathbf R\operatorname{Hom}_R(Y,Z)\) 由 \(\operatorname{Hom}^\bullet_R(Q,I)\) 表示，而且有下界。外层可只替换第一变量，故右侧由
\(\operatorname{Hom}^\bullet_R(P,\operatorname{Hom}^\bullet_R(Q,I))\) 表示。逐次定义
\[
\Phi(f)(p)(q)=f(p\otimes q).
\]
若 \(f\) 的次数为 \(n\)，\(p,q\) 的次数为 \(a,b\)，则值在 \(I^{a+b+n}\)，与两边的次数规定一致。展开右侧的微分得到
\[
\mathrm{d}_I f(p\otimes q)
-(-1)^{a+n}f(p\otimes \mathrm{d}_Qq)
-(-1)^n f(\mathrm{d}_Pp\otimes q),
\]
它正是左侧微分后再施加 \(\Phi\) 的结果。普通张量泛性质逐项给出逆映射；Hom 中的积允许对所有次数同时指定映射，所以这个对应没有有限生成假设。它是自然的复形同构。取零次同调，并用\cref{b4:thm:derived-hom}，得到第二个公式。
\end{proof}

\begin{proposition}\label{b4:prop:derived-extension-restriction}
设 \(f:R\to S\) 为含幺环同态，不要求交换。简记左 \(R\)-模、左 \(S\)-模的导出范畴为 \(D(R),D(S)\)。对有上界的左 \(R\)-模复形 \(X\) 与有下界的左 \(S\)-模复形 \(Y\)，有自然同构
\[
\operatorname{Hom}_{D(S)}(S\otimes_R^{\mathbf L}X,Y)
\simeq\operatorname{Hom}_{D(R)}(X,\operatorname{Res}_R^S Y).
\]
其中导出标量扩张由 \(S\otimes_RP\) 计算，\(P\to X\) 为有上界投射替换，且 \(S\) 带左 \(S\)、右 \(R\) 的通常双模结构。
\end{proposition}
\begin{proof}
\(S\otimes_RP\) 逐项是投射左 \(S\)-模，因为标量扩张把自由模送到自由模并保持直和项。普通伴随逐项给出复形同构
\[
\operatorname{Hom}_S^\bullet(S\otimes_RP,Y)
\simeq\operatorname{Hom}_R^\bullet(P,\operatorname{Res}_R^S Y),
\]
映射为 \(g\mapsto(p\mapsto g(1\otimes p))\)，逆为 \(h\mapsto(s\otimes p\mapsto s\cdot h(p))\)。二者保持 Hom 微分。投射比较把两边的零次同调分别识别为所需导出态射群。
\end{proof}

这一伴随不要求 \(S\) 在 \(R\) 上平坦，因为已经替换 \(X\)。忘却函子本来正合，不另取导出；它与导出标量扩张的关系可对照 Weibel\cite[Theorem 10.7.6]{Weibel1994HomologicalAlgebra}，该书在此使用交换环版本。

\subsection{基变换的充分条件与反例}
\label{b4:subsec:1803:03}

基变换有两种不同的问题。一个是把 \(R\) 换成 \(S\) 后，普通张量是否已能计算导出张量；另一个是张量能否穿过 Hom。前者由平坦性控制，后者还涉及有限生成条件，不能只核对其中一项。

\begin{proposition}\label{b4:prop:flat-base-change-hom}
设 \(R\to S\) 是交换环同态，\(S\) 作为 \(R\)-模平坦。设 \(X\in D^-(R\text{-}\mathbf{Mod})\) 在该导出范畴中同构于有限生成投射 \(R\)-模的有上界复形 \(P\)，\(Y\) 为有下界的 \(R\)-模复形。则自然映射
\[
S\otimes_R\mathbf R\operatorname{Hom}_R(X,Y)
\longrightarrow
\mathbf R\operatorname{Hom}_S(S\otimes_RX,S\otimes_RY)
\]
是同构。左侧及右侧变量中的普通张量均已代表导出标量扩张。
\end{proposition}
\begin{proof}
平坦性使 \(S\otimes_R-\) 保持拟同构，故可用 \(P\) 替换 \(X\)。在这些模型上定义复形映射
\[
\alpha:S\otimes_R\operatorname{Hom}_R^\bullet(P,Y)
\longrightarrow
\operatorname{Hom}_S^\bullet(S\otimes_RP,S\otimes_RY),
\]
其中对次数 \(n\) 的齐次映射 \(f\)、\(x\in P^i\) 及 \(s,s'\in S\)，规定
\[
\alpha(s\otimes f)(s'\otimes x)=ss'\otimes f(x).
\]
\(R,S\) 的交换性保证该公式对张量关系良定义，且所得映射为 \(S\)-线性；将 Hom 微分代入也得到 \(\alpha\delta=\delta\alpha\)。这个公式与两个变量的复形映射相容，因而传到导出范畴后就是命题中的自然映射。

对有限生成投射模 \(P^i\)，它在每个位置给出的映射
\[
S\otimes_R\operatorname{Hom}_R(P^i,Y^j)
\longrightarrow
\operatorname{Hom}_S(S\otimes_RP^i,S\otimes_RY^j)
\]
为同构：有限自由模时是有限个副本的恒等识别，直和项给出投射情形。固定总次数 \(n=j-i\)，有上界的 \(P\) 与有下界的 \(Y\) 只留下有限多个非零位置，所以张量可以与该次数中的 Hom 积交换。以上逐项同构与微分相容，给出复形同构。\(S\otimes_RP\) 逐项投射，故其 Hom 计算右侧。
\end{proof}

\begin{proposition}\label{b4:prop:derived-base-change-perfect}
设 \(R\to S\) 是任意交换环同态，\(X\in D(R\text{-}\mathbf{Mod})\) 在该导出范畴中同构于有限生成投射 \(R\)-模的有界复形，\(Y\in D^-(R\text{-}\mathbf{Mod})\)。则存在自然同构
\[
S\otimes_R^{\mathbf L}\mathbf R\operatorname{Hom}_R(X,Y)
\simeq
\mathbf R\operatorname{Hom}_S
\bigl(S\otimes_R^{\mathbf L}X,S\otimes_R^{\mathbf L}Y\bigr).
\]
\end{proposition}
\begin{proof}
用给定的有界有限生成投射复形 \(P\) 表示 \(X\)，取有上界投射替换 \(Q\to Y\)。由于 \(P\) 有界，\(\operatorname{Hom}_R^\bullet(P,Q)\) 有上界，每项是有限个投射模的直和：\(\operatorname{Hom}_R(P^i,Q^j)\) 是有限个 \(Q^j\) 的直和项。因此它既表示 \(\mathbf R\operatorname{Hom}_R(X,Y)\)，又可直接计算左侧的导出张量。

\(S\otimes_RP\) 仍为有界投射复形，故右侧可由
\(\operatorname{Hom}_S^\bullet(S\otimes_RP,S\otimes_RQ)\) 计算。上一命题中定义的比较映射对有限生成投射模逐项给出 Hom--基变换同构，这一步不需要平坦性；此处每次只有有限个项，把它们合起来便得到自然的复形同构。第二变量无需有下界：\(P\) 与 \(S\otimes_RP\) 均为有界投射复形，\cref{b4:prop:resolution-computes-derived-hom}对任意目标复形的比较结论保证上述 Hom 模型有效，而 \(Q\) 有上界保证左侧的导出张量仍在有上界范围内计算。
\end{proof}

平坦性不能在第一个命题中直接删去。取 \(R=\mathbb Z\)、\(S=\mathbb Z/2\)、\(M=\mathbb Z/2\)、\(N=\mathbb Z\)。则
\[
S\otimes_{\mathbb Z}\operatorname{Hom}_{\mathbb Z}(M,N)=0,
\qquad
\operatorname{Hom}_S(S\otimes M,S\otimes N)=S.
\]
在一次 Ext 上，左侧 \(S\otimes\operatorname{Ext}_{\mathbb Z}^1(M,N)=S\)，右侧 \(\operatorname{Ext}_S^1(S,S)=0\)。这些计算已显示普通基变换不能直接移入 Hom 或 Ext。对前一命题中的全复形比较，用自由消解构造的 Hom 模型，\(\mathbf R\operatorname{Hom}_{\mathbb Z}(S,\mathbb Z)\) 由零、一次的 \(\mathbb Z\xrightarrow{-2}\mathbb Z\) 表示；对这个模型作普通张量，得到两次均为 \(S\)、微分为零的复形，而右侧 \(\mathbf R\operatorname{Hom}_S(S,S)=S[0]\)。失配出现在一次上同调。由于普通张量不保持拟同构，离开平坦假设后，左侧若不指定模型，甚至不是导出范畴上的良定义运算；改成导出张量才能消除这一问题。

\ProseParagraphBreak
有限生成条件也有实质作用。即使 \(\mathbb Q\) 在 \(\mathbb Z\) 上平坦，对 \(P=\bigoplus_{j\ge1}\mathbb Z\)、\(N=\mathbb Z\)，比较映射为
\[
\mathbb Q\otimes_{\mathbb Z}\prod_{j\ge1}\mathbb Z
\longrightarrow\prod_{j\ge1}\mathbb Q.
\]
左侧像中的序列有一个共同分母，右侧的 \((1/j!)_{j\ge1}\) 没有共同分母，所以不满。无限自由虽然投射，仍不能把张量从无限 Hom 积外直接搬入。

\subsection{双数代数上的周期消解与非平坦基变换}
\label{b4:subsec:1803:04}

\begin{example}\label{b4:exm:dual-numbers-derived-base-change}
设 \(k\) 是域，\(R=k[\varepsilon]/(\varepsilon^2)\)，\(k=R/(\varepsilon)\)。取每项为 \(R\) 的有上界自由复形
\[
P=\bigl(\cdots\xrightarrow{\varepsilon}R
\xrightarrow{\varepsilon}R\xrightarrow{\varepsilon}R\bigr),
\qquad P^i=R\ (i\le0),
\]
并用零次的商映射增广到 \(k\)。在 \(R\) 中，乘 \(\varepsilon\) 的核与像都是 \((\varepsilon)\)，故增广正合，\(P\to k\) 是自由替换。与 \(k\) 张量后微分全为零，因此
\[
H^{-n}(k\otimes_R^{\mathbf L}k)=k\quad(n\ge0),\qquad
k\otimes_Rk=k.
\]
普通张量只保留零次项，所有负次同调均被遗漏。这同时说明 \(k\) 在 \(R\) 上不平坦，且投射维数无限；这里是在导出对象上重新观察\cref{b4:exm:infinite-dimension-dual-numbers}的同一个周期消解。

\ProseParagraphBreak
令 \(f:R\to k\)。限制标量本来正合，而再作导出标量扩张会得到 \(k\otimes_R^{\mathbf L}k\)，不是恒等对象 \(k\)。因此“先忘却，再扩张就回到原对象”即使对商环上的自由秩一模也不成立。伴随的余单位确实给出
\(k\otimes_R^{\mathbf L}k\to k\)，在上述代表中它保留零次项、把负次数送到零；它的存在不表示它是同构。

\ProseParagraphBreak
同一个 \(P\) 计算导出 Hom 时，\(\operatorname{Hom}_R(P,k)\) 的微分也全为零，于是
\[
\operatorname{Ext}_R^n(k,k)=k\quad(n\ge0).
\]
张量的额外同调向负次数延伸，Hom 的额外同调向正次数延伸，二者都在记录未能用有限长自由消解代替 \(k\) 的事实。把 \(R\) 直接换成 \(k\) 后，\(k\) 成为自由模，高次 \(\operatorname{Tor}^k\) 和 \(\operatorname{Ext}_k\) 全部消失；这正是普通非平坦基变换丢失的信息。
\end{example}

% !TeX root = ../../Book4.tex
% 纲要标题：### 19.4* 无界情形
% 纲要标签：b4:sec:1804
% 纲要位置：计划纲要.md，第 2971 行。

\OptionalSection{无界情形}{b4:sec:1804}

无界复形没有可供逐度归纳的起点，逐项投射或内射便不再自动足够。K-投射、K-内射与 K-平坦直接要求复形正确处理零调复形，并支持无界替换。\subsecref{b4:subsec:2002:03}中任意余积的构造和紧对象的反向刻画则以这里的半自由替换及其 K-投射性质为前提。

% 纲要内容（仅作写作计划，尚非教材正文）：
%
% * K-内射、K-投射与 K-平坦
% * 无界消解的必要修正
% * 不把有界证明直接无条件推广
%
% ---
%

\subsection{K-内射、K-投射与 K-平坦}
\label{b4:subsec:1804:01}

有上界投射复形和有下界内射复形之所以有效，关键不是每一项的名称，而是整个复形对正合复形的作用。无界时应直接把这种作用列为条件。

\begin{definition}\label{b4:def:k-resolutions}
设 \(R\) 为环，以下均为上链复形。
\begin{enumerate}
\item 左 \(R\)-模复形 \(P\) 若对每个正合左 \(R\)-模复形 \(A\)，都有 \(\operatorname{Hom}_R^\bullet(P,A)\) 正合，则称 \(P\) 为\term[K-toushe]{K-投射复形}[K-projective complex]。
\item 左 \(R\)-模复形 \(I\) 若对每个正合左 \(R\)-模复形 \(A\)，都有 \(\operatorname{Hom}_R^\bullet(A,I)\) 正合，则称 \(I\) 为\term[K-neishe]{K-内射复形}[K-injective complex]。
\item 左 \(R\)-模复形 \(F\) 若对每个正合右 \(R\)-模复形 \(A\)，直和总复形 \(A\otimes_RF\) 都正合，则称 \(F\) 为\term[K-pingtan]{K-平坦复形}[K-flat complex]。右模版本交换两个侧别。
\end{enumerate}
\end{definition}

前两项分别等价于所有 \(P\to A[n]\)、\(A\to I[n]\) 都零同伦，因为 Hom 复形的同调正是这些同伦类。由映射锥判据，这也等价于相应 Hom 函子把拟同构送到拟同构；K-平坦性则等价于张量函子具有这个性质。

\ProseParagraphBreak
\cref{b4:thm:bounded-projective-model}的提升论证说明，有上界的逐项投射复形是 K-投射的；其对偶说明有下界的逐项内射复形是 K-内射的。\cref{b4:lem:bounded-flat-tensor-preserves}说明有上界的逐项平坦复形是 K-平坦的。最后一个结论允许另一个变量无界，原因在于直和总化与滤过余极限的正合性，而不是对角线自动有限。

\begin{example}\label{b4:exm:unbounded-free-not-k-projective}
仍取 \(R=k[\varepsilon]/(\varepsilon^2)\)，令 \(E^n=R\) 对每个整数 \(n\)，所有微分均为乘 \(\varepsilon\)。\(E\) 正合且每项自由，但不是 K-投射：若 \(1_E\) 零同伦，则每度存在乘某元素 \(a_n\) 的同伦，使
\(1=\varepsilon a_n+a_{n+1}\varepsilon\)，右侧属于 \((\varepsilon)\)，矛盾。因此 \(\operatorname{Hom}_R(E,E)\) 的零次同调非零。

\ProseParagraphBreak
它也不是 K-平坦。取\cref{b4:exm:dual-numbers-derived-base-change}中的自由替换 \(P\to k\)。若 \(E\) 是 K-平坦，则
\(E\otimes_RP\to E\otimes_Rk\) 为拟同构。\(P\) 本身有上界且平坦，所以 \(P\otimes_RE\) 正合；交换因子的分次符号给出 \(E\otimes_RP\) 正合。但 \(E\otimes_Rk\) 每度为 \(k\)、微分为零，不正合，矛盾。

\ProseParagraphBreak
这个例子的每项甚至内射。取线性泛函 \(\lambda:R\to k\)，\(\lambda(a+b\varepsilon)=b\)。映射
\(R\to\operatorname{Hom}_k(R,k)\)、\(r\mapsto(s\mapsto\lambda(sr))\)
是 \(R\)-模同构：双线性型在基 \(1,\varepsilon\) 下的矩阵为 \(\left(\begin{smallmatrix}0&1\\1&0\end{smallmatrix}\right)\)。又有
\(\operatorname{Hom}_R(M,\operatorname{Hom}_k(R,k))\simeq\operatorname{Hom}_k(M,k)\)，右侧在向量空间短正合列上正合，故 \(R\) 内射。尽管如此，\(E\) 不是 K-内射，因为 \(E\) 正合而 \(\operatorname{Hom}_R(E,E)\) 不正合。
\end{example}

\subsection{无界消解的修正}
\label{b4:subsec:1804:02}

在模范畴中，任意无界复形仍可得到足够好的替换，但须按复形的全部循环关系构造。下面在模范畴中构造这样的替换。

\begin{definition}\label{b4:def:semifree-complex}
设 \(R\) 为环，\(P\) 是左 \(R\)-模复形。若存在子复形滤过
\[
0=P_{-1}\subseteq P_0\subseteq P_1\subseteq\cdots,
\qquad P=\bigcup_{s\ge0}P_s,
\]
使每个包含在分次模意义下分裂，每个商 \(P_s/P_{s-1}\) 是若干移位 \(R[n]\) 的直和且微分为零，则称此 \(P\) 连同滤过为\term[banziyou-fuxing]{半自由复形}[semi-free complex]。等价地，每一层可添一批齐次自由生成元，新生成元的微分只涉及更早层的生成元。
\end{definition}

\begin{theorem}\label{b4:thm:semifree-resolution}
设 \(R\) 为环，\(X\) 为任意左 \(R\)-模上链复形。存在半自由复形 \(P\) 及拟同构 \(P\to X\)。每个半自由复形都是 K-投射且 K-平坦的；这里不要求 \(X\) 或 \(P\) 有任何次数界。
\end{theorem}
\begin{proof}
先对 \(X\) 每个次数的每个余圈 \(z\in Z^nX\)，添一个次数 \(n\) 的自由生成元 \(e_z\)，令其微分为零，映到 \(z\)。这些余圈由集合 \(\coprod_{n\in\mathbb Z}Z^nX\) 索引，所得 \(P_0\to X\) 在同调上满射。

假定已构造 \(f_s:P_s\to X\)。对每对数据
\[
z\in Z^nP_s,\qquad x\in X^{n-1},\qquad f_s(z)=\mathrm{d}_Xx,
\]
添一个次数 \(n-1\) 的自由生成元 \(e_{z,x}\)；本层全部数据是集合 \(\coprod_{n\in\mathbb Z}(Z^nP_s\times X^{n-1})\) 的子集。规定
\(\mathrm{d} e_{z,x}=z\)、\(f_{s+1}(e_{z,x})=x\)。因为 \(\mathrm{d}z=0\) 且 \(f_s(z)=\mathrm{d}x\)，复形条件与映射条件同时成立。这定义 \(P_{s+1}\)，其分次模是 \(P_s\) 与所添自由模的直和，商微分为零。取并得到半自由复形 \(P\) 及映射 \(f:P\to X\)。

每个 \(X\) 的同调类在 \(P_0\) 已有余圈代表，因而 \(H(f)\) 满。若 \(z\in Z^nP\) 映成 \(X\) 中的边界，则 \(z\) 只含有限个生成元，属于某个 \(P_s\)；选定 \(x\) 使 \(f(z)=\mathrm{d}x\)，下一层就添入 \(e_{z,x}\) 使 \(z\) 成为边界。因此 \(H(f)\) 单，\(f\) 是拟同构。

为证 K-投射，设 \(A\) 正合，\(g:P\to A\) 是复形映射。逐层构造次数 \(-1\) 的映射 \(h\)，使 \(g=\mathrm{d}h+h\mathrm{d}\)。若 \(h\) 已在旧层定义，对次数 \(n\) 的新生成元 \(e\)，\(\mathrm{d}e\) 属于旧层，而且
\[
\mathrm{d}_A\bigl(g(e)-h(\mathrm{d}e)\bigr)=g(\mathrm{d}e)-\mathrm{d}_Ah(\mathrm{d}e)=0,
\]
后一等式用旧层上的同伦公式及 \(\mathrm{d}^2e=0\)。由 \(A\) 正合，每个 \(g(e)-h(\mathrm{d}e)\) 都有原像。按\subsecref{b4:subsec:0101:03}的集合指标选择约定，可为本层全部新生成元同时选取 \(h(e)\in A^{n-1}\)，使其微分为 \(g(e)-h(\mathrm{d}e)\)。自由性允许线性延拓到本层。逐层取并即得全局同伦。把 \(A\) 换成任意移位，得到所有 Hom 同调消失。

为证 K-平坦，设 \(A\) 是正合右模复形。每个商
\(A\otimes_R(P_s/P_{s-1})\) 是 \(A\) 的移位直和，所以正合。滤过包含逐次分裂，张量后保持短正合列，归纳得 \(A\otimes_RP_s\) 正合。最后由滤过余极限的正合性，\(A\otimes_RP\) 正合。
\end{proof}

这里的归纳不是沿复形次数开始，而是沿添加生成元的层次开始。即使次数向两边无限延伸，每个新生成元的微分只依赖已经处理的部分，所以构造同伦仍有起点。

\subsection{有界论证向无界情形推广的条件}
\label{b4:subsec:1804:03}

\begin{proposition}\label{b4:prop:k-resolution-comparison}
设 \(R\) 为环，简记左 \(R\)-模的同伦范畴和导出范畴为 \(K(R)=K(R\text{-}\mathbf{Mod})\)、\(D(R)=D(R\text{-}\mathbf{Mod})\)。若 \(P\) 为 K-投射左 \(R\)-模上链复形，则对任意左 \(R\)-模上链复形 \(Y\)，有
\[
\operatorname{Hom}_{K(R)}(P,Y)\simeq\operatorname{Hom}_{D(R)}(P,Y).
\]
若 \(I\) 为 K-内射左 \(R\)-模上链复形，则对任意左 \(R\)-模上链复形 \(X\)，有对偶同构
\(\operatorname{Hom}_{K(R)}(X,I)\simeq\operatorname{Hom}_{D(R)}(X,I)\)。
\end{proposition}
\begin{proof}
设 \(s:Y'\to Y\) 为拟同构。\(\operatorname{Cone}(s)\) 正合，K-投射性及 Hom 的锥长正合列说明
\(\operatorname{Hom}_{K(R)}(P,Y')\to\operatorname{Hom}_{K(R)}(P,Y)\)
是同构。因此从 \(P\) 出发的同伦态射函子已经把所有分母变成同构。

更具体地，导出态射的屋顶 \(P\xleftarrow{s}U\xrightarrow{a}Y\) 中，\(s\) 对 \(\operatorname{Hom}_{K(R)}(P,-)\) 给出同构，故唯一有 \(t:P\to U\) 使 \(st=1_P\) 在同伦范畴中成立。由拟同构的二取三性质，\(t\) 也是拟同构。因此原屋顶与 \(P\xleftarrow{1_P}P\xrightarrow{at}Y\) 以 \(P\) 为共同细化对象，细化映射分别为 \(t:P\to U\) 与 \(1_P:P\to P\)：左侧由 \(st=1_P\) 交换，右侧由 \(a\circ t=at\) 交换。故原屋顶由普通态射 \(at\) 表示。若两个这样的态射在导出范畴中相等，分式等价关系给出一个共同分母；对该分母同样施加上述 Hom 同构，便得它们已在同伦范畴中相等。对偶证明使用反向屋顶及 \(\operatorname{Hom}_{K(R)}(-,I)\)。
\end{proof}

右 \(R\)-模等同于左 \(R^{\mathrm{op}}\)-模，所以\cref{b4:def:semifree-complex,b4:thm:semifree-resolution}应用于反环 \(R^{\mathrm{op}}\) 时，给出右半自由复形及其替换；滤过商是移位右正则模 \(R_R[n]\) 的直和。给定任意无界右 \(R\)-模复形 \(X\) 与左 \(R\)-模复形 \(Y\)，分别取右半自由替换 \(P\to X\) 与左半自由替换 \(Q\to Y\)。直和总复形 \(P\otimes_RQ\) 表示 \(X\otimes_R^{\mathbf L}Y\)。K-平坦性使 \(P\otimes_RY\) 或 \(X\otimes_RQ\) 也能计算同一导出对象，K-投射比较保证态射和不同替换之间的相容性。

\ProseParagraphBreak
对任意无界左 \(R\)-模复形 \(X,Y\)，导出 Hom 可在第一变量用左半自由替换 \(P\to X\)，定义 \(\mathbf R\operatorname{Hom}_R(X,Y)=\operatorname{Hom}_R^\bullet(P,Y)\)，其中 Hom 仍取各次数的直积。无界替换的作用来自 K-投射、K-平坦性质，直和总化的对角线上可以有无限多个非零项。

\ProseParagraphBreak
对一般交换范畴，未必有自由生成元，也未必有所需的无限直和或精确余极限。若要用 K-内射复形构造无界 \(RF\)，必须另证每个对象有相应替换，并检查 \(F\) 在这些替换间保持拟同构；仅有“足够内射对象”并不自动提供这里使用的整个复形性质。可比较 Weibel\cite[Generalized Existence Theorem 10.5.9]{Weibel1994HomologicalAlgebra}所列出的替换与零调性条件；有界性有限制时，前面各节已经给出了可直接使用的构造。


\begin{chaptersummary}
全导出运算先选择能控制拟同构的复形，再进行原有加性运算。内射模型给出全右导出，投射模型给出全左导出；各项对相应函子零调的单侧有界复形也能计算相应结果。复合导出之间总有比较，只有满足相应零调性条件时才得到同构。

\ProseParagraphBreak
导出张量的负次同调是 Tor，导出 Hom 的正次同调在模上是 Ext。它们之间的伴随来自带分次符号的复形同构。平坦基变换和有限投射复形分别提供两类可验证的基变换条件；双数代数说明非平坦的普通张量可能同时丢掉无限多项。无界情形的修正，则是用 K-性质控制整个复形，而不只检查各项。
\end{chaptersummary}

\BookExercises

\begin{exercise}\label{b4:ex:18-additive-right-derived}
取加性函子 \(F:\mathbf{Ab}\to\mathbf{Ab}\)，\(F(M)=\operatorname{Hom}_{\mathbb Z}(\mathbb Z/2\mathbb Z,M)\)。证明 \(LF:D^-(\mathbf{Ab})\to D^-(\mathbf{Ab})\) 是零函子，并计算 \(RF(\mathbb Z)\)。分别判定比较态射 \(F(\mathbb Z)[0]\to RF(\mathbb Z)\) 与 \(LF(\mathbb Z/2\mathbb Z)\to F(\mathbb Z/2\mathbb Z)[0]\) 是否为同构，说明两种导出构造与 \(F\) 的正合方向有何关系。
\end{exercise}
\begin{solution}
投射交换群是自由交换群的直和因子，因而无挠；从 \(\mathbb Z/2\mathbb Z\) 到它的同态只能为零。所以任意有上界投射替换逐项施加 \(F\) 后都为零，\(LF=0\)。另一方面，用\cref{b4:exm:integer-injective-resolution}的内射消解
\[
0\longrightarrow\mathbb Z\longrightarrow\mathbb Q\longrightarrow\mathbb Q/\mathbb Z\longrightarrow0
\]
计算全右导出。\(F(\mathbb Q)=0\)，而 \(F(\mathbb Q/\mathbb Z)\) 是由 \(\dfrac{1}{2}+\mathbb Z\) 生成的二阶子群。因此 \(RF(\mathbb Z)\simeq(\mathbb Z/2\mathbb Z)[-1]\)，其唯一非零同调在一次。

\(F(\mathbb Z)=0\)，故第一个比较是零对象到非零对象的态射，不是同构。\(F(\mathbb Z/2\mathbb Z)\simeq\mathbb Z/2\mathbb Z\)，而 \(LF(\mathbb Z/2\mathbb Z)=0\)，第二个比较也不是同构。\(F\) 左正合，所以模上的 \(H^0RF\) 恢复 \(F\)，但更高次仍可能非零。它不右正合：满射 \(\mathbb Z\to\mathbb Z/2\mathbb Z\) 被送到 \(0\to\mathbb Z/2\mathbb Z\)，不再满射，因而不能期待 \(H^0LF\) 恢复 \(F\)。这与\cref{b4:prop:additive-right-derived-zero}中右正合函子的现象相对照。
\end{solution}

\begin{exercise}\label{b4:ex:18-derived-annihilation}
固定正整数 \(m\)，对任意交换群 \(N\)，用两项复形表示 \(\mathbf R\operatorname{Hom}_{\mathbb Z}(\mathbb Z/m\mathbb Z,N)\)。证明在这个导出对象上乘 \(m\) 的态射为零，并给出同伦。
\end{exercise}
\begin{solution}
自由替换 \(\mathbb Z\xrightarrow{m}\mathbb Z\) 位于 \(-1,0\) 次。其 Hom 复形在 \(0,1\) 次为 \(N\xrightarrow{-m}N\)；把一次项乘 \(-1\)，可写成 \(C=(N\xrightarrow{m}N)\)。取唯一可能的同伦分量 \(h^1:C^1\to C^0\) 为恒等，其余为零，则 \(\mathrm{d}h+h\mathrm{d}=m1_C\)。所以乘 \(m\) 在同伦范畴中已为零，而不只是对每个 Ext 群诱导零映射。
\end{solution}

\begin{exercise}\label{b4:ex:18-integer-derived-tensor}
设 \(m,n>0\)，\(d=\gcd(m,n)\)。求 \(\mathbb Z/m\mathbb Z\otimes_{\mathbb Z}^{\mathbf L}\mathbb Z/n\mathbb Z\) 的一个复形代表及全部同调，并判定它何时为零对象。
\end{exercise}
\begin{solution}
代表为 \(\mathbb Z/n\mathbb Z\xrightarrow{m}\mathbb Z/n\mathbb Z\)，次数 \(-1,0\)。乘 \(m\) 的核由 \(n/d\) 的类生成，阶为 \(d\)，余核也为 \(\mathbb Z/d\mathbb Z\)。因此 \(H^{-1}=H^0=\mathbb Z/d\mathbb Z\)，其余为零。导出对象为零当且仅当 \(d=1\)。当 \(d>1\) 时，仅列出普通张量 \(\mathbb Z/d\mathbb Z\) 不足以表示该导出对象。
\end{solution}

\begin{exercise}\label{b4:ex:18-tensor-shift-sign}
设 \(R\) 为环，\(C\) 为右 \(R\)-模复形，\(D\) 为左 \(R\)-模复形，\(a,b\in\mathbb Z\)。构造复形同构
\(C[a]\otimes_RD[b]\simeq(C\otimes_RD)[a+b]\)，并核对符号。
\end{exercise}
\begin{solution}
若 \(x\in C[a]^p=C^{p+a}\)、\(y\in D[b]^q=D^{q+b}\)，定义 \(x\otimes y\mapsto(-1)^{bp}x\otimes y\)。两边总次数都为 \(p+q\)。先在左边微分再映射，两个系数分别为 \((-1)^{a+b(p+1)}\) 与 \((-1)^{p+b+bp}\)。先映射再在右边微分，系数分别为 \((-1)^{bp+a+b}\) 与 \((-1)^{bp+a+b+p+a}\)，逐项相同。再乘相同的符号即为逆。采用投射或平坦替换后得到相同的导出移位同构。
\end{solution}

\begin{exercise}\label{b4:ex:18-extension-tower}
设 \(R\to S\to T\) 为含幺环同态，\(X\) 为有上界的左 \(R\)-模复形。证明
\[
T\otimes_S^{\mathbf L}(S\otimes_R^{\mathbf L}X)
\simeq T\otimes_R^{\mathbf L}X,
\]
其中每次均按标量扩张保留左侧环的作用，不假定平坦性。
\end{exercise}
\begin{solution}
取 \(P\to X\) 的有上界投射替换。\(S\otimes_RP\) 逐项投射于 \(S\)，所以外层导出张量无需再次替换。左侧由 \(T\otimes_S(S\otimes_RP)\) 表示，普通结合映射
\(t\otimes(s\otimes p)\mapsto ts\otimes p\)
给出到 \(T\otimes_RP\) 的复形同构，逆为 \(t\otimes p\mapsto t\otimes(1\otimes p)\)。两边已是各自导出计算，故无须要求 \(S\) 或 \(T\) 平坦。
\end{solution}

\begin{exercise}\label{b4:ex:18-integer-derived-hom}
设 \(m,n>0\)。计算 \(\mathbf R\operatorname{Hom}_{\mathbb Z}(\mathbb Z/m\mathbb Z,\mathbb Z/n\mathbb Z)\)，并与上一道整数导出张量题比较非零次数。
\end{exercise}
\begin{solution}
自由替换给出 \(\mathbb Z/n\mathbb Z\xrightarrow{-m}\mathbb Z/n\mathbb Z\)，次数 \(0,1\)。故零次为乘 \(m\) 的核，一次为余核，均同构于 \(\mathbb Z/\gcd(m,n)\mathbb Z\)，其他次数为零。张量的修正项位于 \(-1\) 次，Hom 的修正项位于 \(1\) 次；这来自 Hom 的第一变量反变，而不是两者采用不同的微分升降约定。
\end{solution}

\begin{exercise}\label{b4:ex:18-finite-projective-bidual}
设 \(R\) 为交换环，\(P\) 为有限生成投射 \(R\)-模的有界复形，\(P^\vee=\operatorname{Hom}_R^\bullet(P,R)\)。证明带符号的求值给出 \(P\simeq(P^\vee)^\vee\)，并解释相应的导出双对偶同构。
\end{exercise}
\begin{solution}
对 \(x\in P^p\)、\(f\in(P^\vee)^{-p}\)，定义 \(\eta(x)(f)=(-1)^pf(x)\)。因为底复形 \(R\) 只有零次，\(\mathrm{d}_{P^\vee}f=-(-1)^{|f|}f\mathrm{d}_P\)。对次数 \(-p-1\) 的 \(g\)，代入两次对偶微分得到
\((\mathrm{d}\eta(x))(g)=(-1)^{p+1}g(\mathrm{d}_Px)=\eta(\mathrm{d}_Px)(g)\)。所以 \(\eta\) 是复形映射。各项有限生成投射，通常的求值到双对偶为同构，乘符号不改变可逆性。\(P\) 及 \(P^\vee\) 都有界投射，所以可直接计算两次导出 Hom，得到 \(P\simeq\mathbf R\operatorname{Hom}_R(\mathbf R\operatorname{Hom}_R(P,R),R)\)。
\end{solution}

\begin{exercise}\label{b4:ex:18-tensor-cone}
设 \(R\) 为交换环，\(f:P\to Q\) 是有上界投射复形间的映射，\(U\) 为有上界投射复形。证明 \(\operatorname{Cone}(f)\otimes_RU\) 自然同构于 \(\operatorname{Cone}(f\otimes1_U)\)，并推出固定一个变量的导出张量保持好三角。
\end{exercise}
\begin{solution}
按分量识别两边为 \((Q\otimes U)\oplus(P\otimes U)[1]\)。若 \(x\in P^{p+1}\)、\(u\in U^q\)，源中把 \(x\) 看作 \(P[1]^p\)，其内部微分为 \(-\mathrm{d}_Px\otimes u+(-1)^px\otimes \mathrm{d}_Uu\)，恰为目标移位总微分。跨分量映射均为 \(f(x)\otimes u\)，故该直接识别就是复形同构。导出范畴每个好三角同构于锥三角；先在投射模型中表示，再作上述张量，所得仍为锥三角。移位的相容性由上一移位题的公式核对。
\end{solution}

\begin{exercise}\label{b4:ex:18-restriction-not-full}
设 \(S=\mathbb Z/2\mathbb Z\)。对商环映射 \(\mathbb Z\twoheadrightarrow S\)，通常模范畴中的限制函子 \(S\text{-}\operatorname{Mod}\to\mathbb Z\text{-}\operatorname{Mod}\) 全忠实。证明导出限制函子在有界导出范畴中未必全忠实。
\end{exercise}
\begin{solution}
商环映射满射，所以两个 \(S\)-模之间的每个 \(\mathbb Z\)-线性映射自动保持 \(S\)-作用，通常的限制函子因而全忠实。再取两个导出对象 \(S\) 和 \(S[1]\)。在 \(D(S)\) 中，\(\operatorname{Hom}(S,S[1])=\operatorname{Ext}_S^1(S,S)=0\)，因为 \(S\) 是自由 \(S\)-模。在 \(D(\mathbb Z)\) 中，同一态射群为 \(\operatorname{Ext}_{\mathbb Z}^1(S,S)=S\ne0\)。其中的非零类由短正合列
\[
0\longrightarrow\mathbb Z/2\mathbb Z
\xrightarrow{\ \bar1\mapsto\bar2\ }\mathbb Z/4\mathbb Z
\longrightarrow\mathbb Z/2\mathbb Z\longrightarrow0
\]
表示，最后一个映射是模 \(2\)。它不分裂，因为 \(\bar1\) 的每个提升都是四阶元，不可能是二阶元的像；中间项也不是 \(S\)-模，因为乘 \(2\) 在其中不为零。有界嵌入全忠实，故这两个计算也分别是 \(D^b(S)\)、\(D^b(\mathbb Z)\) 中的态射群。因此导出限制函子在这对对象上不满。零次模同态没有改变，但允许扩张以后，底环提供了新的态射。
\end{solution}

\begin{exercise}\label{b4:ex:18-bounded-not-tensor-closed}
设 \(k\) 为域，\(R=k[\varepsilon]/(\varepsilon^2)\)。证明 \(D^b(R\text{-}\operatorname{Mod})\) 对导出张量未必封闭。再给出保证 \(X\otimes_R^{\mathbf L}Y\) 有界的一项充分条件，其中 \(X,Y\) 为 \(R\)-模上链复形，\(Y\) 有界。
\end{exercise}
\begin{solution}
取 \(X=Y=k\)，二者只有零次同调，故属于有界导出范畴。但\cref{b4:exm:dual-numbers-derived-base-change}给出每个 \(n\ge0\) 的 \(H^{-n}(k\otimes_R^{\mathbf L}k)=k\)，结果不有界。一个充分条件是 \(X\) 同构于有界平坦复形 \(F\)：此时 \(F\otimes_RY\) 计算导出张量，两个有界次数区间之和仍是有限区间，所以结果有界。
\end{solution}

\begin{exercise}\label{b4:ex:18-cohomology-does-not-determine}
设 \(k\) 为域，\(R=k[\varepsilon]/(\varepsilon^2)\)，\(C=(R\xrightarrow{\varepsilon}R)\) 位于 \(-1,0\) 次，\(T=k[1]\oplus k\)，其中 \(k=R/(\varepsilon)\)。对每个整数 \(n\)，计算 \(\operatorname{Hom}_{D(R)}(C,k[n])\) 与 \(\operatorname{Hom}_{D(R)}(T,k[n])\) 这两个 \(k\)-向量空间，并用结果区分 \(C\) 与 \(T\)。
\end{exercise}
\begin{solution}
\(C\) 有界自由，可以直接用 \(\operatorname{Hom}_R^\bullet(C,k)\) 计算。这是位于 \(0,1\) 次的两项 \(k\)，微分为零，因为 \(\varepsilon\) 在 \(k\) 上作用为零。所以
\[
\operatorname{Hom}_{D(R)}(C,k[n])\simeq
\begin{cases}k,&n=0,1,\\0,&\text{其他 }n.\end{cases}
\]
由\cref{b4:exm:dual-numbers-derived-base-change}中的周期自由消解，\(\operatorname{Ext}_R^j(k,k)=k\) 对每个 \(j\ge0\) 成立；负次导出 Hom 为零。移位第一变量给出
\[
\operatorname{Hom}_{D(R)}(T,k[n])
\simeq\operatorname{Ext}_R^{n-1}(k,k)\oplus\operatorname{Ext}_R^n(k,k)
\simeq\begin{cases}0,&n<0,\\k,&n=0,\\k\oplus k,&n\ge1.\end{cases}
\]
在 \(n=1\) 时两者的维数不同，故 \(C\not\simeq T\)。这用移位态射群重新检验了\cref{b4:prop:derived-cohomology-not-split}：这里同调仅在两度非零，问题是这些同调之间的扩张信息；前一题则是两个有界对象的导出张量产生了无限多度同调。
\end{solution}

\begin{exercise}\label{b4:ex:18-rational-vanishing}
设 \(M\) 为交换群。证明 \(\mathbb Q\otimes_{\mathbb Z}^{\mathbf L}M\) 为零对象，当且仅当 \(M\) 的每个元素都有有限阶。
\end{exercise}
\begin{solution}
\(\mathbb Q\) 是 \(\mathbb Z\) 关于非零整数的局部化，所以平坦，导出张量就是零次的 \(\mathbb Q\otimes M\)。局部化的零判据说明 \(1\otimes m=0\) 当且仅当存在非零整数 \(s\) 使 \(sm=0\)。这些元素生成张量积，故它为零当且仅当所有 \(m\) 都为挠元。此处不要求存在同时零化整个群的同一个整数。
\end{solution}

\begin{optionalexercise}\label{b4:ex:18-acyclic-k-projective}
证明正合的 K-投射复形可缩，正合的 K-内射复形也可缩。反过来，证明每个可缩模复形都同时具有 K-投射、K-内射与适当侧别的 K-平坦性。
\end{optionalexercise}
\begin{solution}
若 \(P\) 正合且 K-投射，在定义中取 \(A=P\)，则 \(\operatorname{Hom}_R^\bullet(P,P)\) 正合。其零次余圈 \(1_P\) 因而是余边，即存在次数 \(-1\) 的 \(h\) 使
\[
1_P=\mathrm{d}_{\operatorname{Hom}}h=\mathrm{d}_Ph+h\mathrm{d}_P.
\]
这就是收缩同伦。K-内射时在定义中取 \(A=I\)，同理得到 \(1_I=\mathrm{d}_Ih+h\mathrm{d}_I\)。反向若 \(1_C=\mathrm{d}h+h\mathrm{d}\)，与任意映射复合即可使从 \(C\) 出发或到 \(C\) 的复形映射零同伦，对各次移位同样成立，所以相应 Hom 复形正合。对右模复形 \(A\)，\(A\otimes C\) 的收缩为 \(a\otimes c\mapsto(-1)^{|a|}a\otimes h(c)\)；代入总微分，两个混合项相消，剩下恒等，故 K-平坦。
\end{solution}

\begin{optionalexercise}\label{b4:ex:18-acyclic-k-flat}
设 \(R\) 为交换环，\(F\) 是正合且 K-平坦的复形。证明对任意模复形 \(Y\)，\(F\otimes_RY\) 都正合。为什么仅假定 \(F\) 逐项平坦不够？
\end{optionalexercise}
\begin{solution}
取半自由替换 \(Q\to Y\)。\(Q\) 是 K-平坦，而 \(F\) 正合，故 \(F\otimes_RQ\) 正合。另一方面，\(F\) 的 K-平坦性使 \(F\otimes_RQ\to F\otimes_RY\) 为拟同构，故目标正合。对\cref{b4:exm:unbounded-free-not-k-projective}的双向周期复形 \(E\)，各项自由且 \(E\) 正合，但 \(E\otimes_Rk\) 零微分、每度非零，因此逐项平坦并不足够。
\end{solution}

\begin{optionalexercise}\label{b4:ex:18-field-unbounded}
设 \(k\) 是域。证明任意无界 \(k\)-向量空间复形都同时是 K-投射、K-内射和 K-平坦的，并说明这一结论为什么不能推广到任意环。
\end{optionalexercise}
\begin{solution}
对任意正合向量空间复形 \(A\)，逐次选补空间 \(A^n=Z^nA\oplus L^n\)。微分把 \(L^n\) 同构到 \(Z^{n+1}A\)，其逆在 \(Z^{n+1}A\) 上定义收缩，补空间上取零，所以 \(A\) 可缩。任意复形 \(X\) 与这个收缩复合，表明 \(\operatorname{Hom}(X,A)\)、\(\operatorname{Hom}(A,X)\) 可缩；张量后也由带分次符号的同伦可缩。因而三个 K-条件都成立。一般环上的正合列未必分裂，正合复形未必可缩；双数代数的双向周期复形已经给出失败例子。
\end{solution}
